
Calibration of neural network confidence scores has been extensively studied for image classification, where temperature scaling achieves 5-15\% ECE reduction \citep{guo2017calibration}. However, calibration for code generation tasks remains largely unquantified. Code generation differs from classification in three ways: (1) autoregressive probability aggregation across tokens, (2) binary evaluation criteria (code passes all tests or fails), and (3) training objectives decoupled from functional correctness. These differences suggest that calibration characteristics may differ from classification tasks.

This work addresses a technical gap: establishing a validated experimental protocol for measuring calibration quality in code generation. The primary contribution is not empirical findings about model calibration, but rather validation of the experimental pipeline needed to obtain such findings. Specifically, this work validates that:

\begin{enumerate}
\item Temperature scaling optimization converges reliably for code generation calibration data
\item ECE computation with 15-bin partitioning produces stable measurements
\item Data splitting strategies (calibration vs. validation) are implementable for MBPP
\item The complete experimental workflow executes without technical failures
\end{enumerate}

\subsubsection{Simulation Mode Rationale}

This work uses simulation mode rather than actual Code Llama 7B inference for three reasons:

\textbf{Technical validation focus:} The goal is to validate experimental infrastructure (optimization convergence, metric computation, visualization generation) independently from model-specific results. Simulation enables rapid iteration on the experimental protocol.

\textbf{Resource efficiency:} Full Code Llama 7B experiments require 4-6 hours on A100 GPUs for 395 problems. Simulation validates the pipeline in minutes, enabling verification that all components function before committing computational resources.

\textbf{Explicit scope:} By using simulation, the work clearly delimits what is validated (experimental protocol) from what requires future work (empirical model behavior). This prevents premature claims about model calibration before production experiments are run.

The simulation uses binary logit representation (correct/incorrect classes) with high baseline ECE (0.527) to reflect the overconfidence patterns observed in prior work on language models \citep{kadavath2022language}. While temperature values from simulation (T* = 2512.71) are artifacts of the binary representation, the core technical findings (LBFGS convergence, ECE reduction mechanism, accuracy preservation) are implementation-independent.

\subsubsection{Relationship to Broader Research Goals}

This technical validation is part of a sequential hypothesis testing protocol for confidence-based iteration control in agentic code generation. The complete hypothesis proposes that calibrated confidence scores enable adaptive resource allocation: high-confidence predictions submitted directly, low-confidence predictions routed to execution feedback. However, this broader hypothesis requires four sequential validations:

\begin{itemize}
\item \textbf{H-E1 (this work):} Temperature scaling produces calibrated confidence scores
\item \textbf{H-M1 (future):} Calibrated confidence correlates monotonically with correctness
\item \textbf{H-M2 (future):} Self-critique benefit decreases with confidence
\item \textbf{H-C1 (future):} Confidence-based gating reduces execution attempts
\end{itemize}

This work completes H-E1 in simulation mode. The remaining hypotheses are not addressed.
